% Reviewed-record descendant: method and scope before corpus results.
\section{Diagnostic Method and Objective Scope}
\label{sec:diagnostic-method}

For a prompt with a sampled group of rewards, the submitted runner centers each
reward by the group mean. If all rewards are homogeneous, every centered
reward-contrast coefficient is zero; hence the submitted reward-contrast term
is zero. We call the fraction of such groups ZVF and write its submitted
complement as Gradient Utilization, $\mathrm{GU}=1-\mathrm{ZVF}$. This is an
accounting identity for this runner, not a measurement of a total optimizer gradient.

The diagnostic runner is \emph{not canonical GRPO}. It was deliberately minimal
for a narrow audit of centered rewards: it omits a PPO probability ratio and
clipping, frozen-reference KL, and completion-only token masking. Those omissions
do not make the omitted terms irrelevant. Clipping and completion-only masking can
change which token contributions survive, while reference KL or auxiliary losses
can make the total gradient nonzero even when the centered reward-contrast term is
zero. The results therefore neither establish canonical-GRPO equivalence nor
transfer to a different objective.

We organize observations by evidence tier and estimand rather than pooling
heterogeneous models, tasks, runners, checkpoints, and evaluators as replications.
In the reviewed record, the two synthetic tool-use cells have reward 0, ZVF 1,
GU 0, and held-out outcome \emph{not evaluated}; GSM8K ZVF is only a heterogeneous
descriptive aggregate unless a named per-run source is given. Later files such as
\texttt{zvf\_iter70\_per\_run\_zvf\_acc.tsv} are future-paper evidence and do not
repair the reviewed record. Under binary rewards, high-reward/high-ZVF is a
theoretical solved-task regime, not an observed claim here: it yields zero centered
contrast, may waste rollout budget, and cannot by unsigned ZVF alone choose
retire/distill over retry/resample/warm-start.
